\documentclass[10pt]{article}

\usepackage[T1]{fontenc}

\usepackage[letterpaper,margin=1in]{geometry}

\usepackage{amsmath}

\usepackage[hidelinks]{hyperref}

\setlength{\parindent}{0pt}

\setlength{\parskip}{.45em}

\emergencystretch=3em

\begin{document}

\section*{Response to the supplied deep review}

Artifact: PDP2-CR030-R1. Local revision prepared on 2 November 2026.

This response covers all 58 detailed review rows, the eight acceptance groups and all 22 checklist items. The review is valuable, but missing supplementary files in its input must be distinguished from missing proofs: the previous package already contained S1--S14. The new delivery makes the supplement visible as a standalone PDF and adds S15.

No channel generation, optimizer campaign or training was performed for this revision. Stored scientific results and all 3,627 reproduction payloads are unchanged. Bounded analytic checks and saved-data statistical replay are separate validation operations. Pricing still adds no achieved-rate gain over corresponding unpriced search and no main/\allowbreak scale 0.01-nat success. The all-zero original ASU failure and the separate correction/\allowbreak postprocessing deviation remain disclosed.

The article retains the 13-page initial-submission layout; the supplement has 32 pages. The final delivery includes the revised article/\allowbreak source ZIP, supplement/\allowbreak source/\allowbreak code/\allowbreak results ZIP, selected-reference exports and cover letter. Public release and author-controlled submission declarations remain separate. Historical generator VCS/\allowbreak scenario tags, BLAS build and separate HiGHS version were not captured and are not invented; installed-source/\allowbreak input hashes and recorded runtime metadata are supplied.

The initial planning audit failed solely because its read-only boundary was ineffective. The protected second planning attempt passed three reflections; the first failure remains recorded. Scientific/\allowbreak manuscript acceptance is reported separately in the independent final review.

\subsection*{R01: Abstract grammar}

The abstract now uses a complete grammatical statement of fixed-power scheduling hardness. The precise NP-completeness statement first appears after the decision-language definition.

{\small Evidence: abstract.tex, model-analysis.tex, IDEAL-MMSE-PAIR.}

\subsection*{R02: Positive powers}

The abstract and introduction specify strictly positive bounded transmit powers.

{\small Evidence: abstract.tex, introduction.tex.}

\subsection*{R03: Seventeen candidates}

The abstract states that rectangular pricing requires at most seventeen candidate points; the theorem and supplement distinguish this characterization from precision-dependent certified computation.

{\small Evidence: abstract.tex, model-analysis.tex Theorem3, full-proofs.tex D6--D8.}

\subsection*{R04: Negative empirical result}

The abstract explicitly states no achieved-rate improvement over unpriced search and no 0.01-nat target success. The adverse result is unchanged.

{\small Evidence: abstract.tex, evaluation-results.tex.}

\subsection*{R05: Exactly-once service}

The introduction now says serving every user exactly once.

{\small Evidence: introduction.tex.}

\subsection*{R06: Article roadmap}

A concise roadmap identifies Sections II--VIII and the directly relevant supplementary sections.

{\small Evidence: introduction.tex, final paragraph.}

\subsection*{R07: Supplement availability}

The original delivery already included the full supplement. The new folder additionally provides a standalone 32-page Supplementary\_\allowbreak Material.pdf, and the ZIP contains that PDF plus every source. Missing material in the review input is distinguished from a missing proof.

{\small Evidence: Supplementary\_\allowbreak Material.pdf, Supplementary\_\allowbreak Materials.zip/\allowbreak source.}

\subsection*{R08: Transmitted symbols}

The signal model explicitly defines the complex Gaussian symbol distribution, unit variance and independence.

{\small Evidence: model-analysis.tex, received signal.}

\subsection*{R09: Whitening wording}

The scalar-noise explanation attributes channel division to whitening and excludes reapplication of path loss.

{\small Evidence: model-analysis.tex, received signal.}

\subsection*{R10: Normalized received power}

The zero-gain convention now uses normalized received power, distinguishing it from transmitted watts.

{\small Evidence: model-analysis.tex, channel summaries.}

\subsection*{R11: Generic correlation}

The generic pair-rate function is introduced with x,y>=0 and r in [0,1].

{\small Evidence: model-analysis.tex, pair log-rate.}

\subsection*{R12: Log-rate units}

The generic rate is stated in nats per channel use.

{\small Evidence: model-analysis.tex, pair log-rate.}

\subsection*{R13: Selected pair}

The selected pair e\_\allowbreak s is defined before the aggregate SE equation.

{\small Evidence: model-analysis.tex, sum-SE definition.}

\subsection*{R14: Aggregate and occupied-band SE}

Q and Q/\allowbreak S are distinguished with their units and resource-bandwidth normalization. The supplement also distinguishes stored tones from actually occupied tones.

{\small Evidence: model-analysis.tex, sum-SE definition, revision-details.tex R1.}

\subsection*{R15: Binary domain}

The binary domain was already explicit in the sentence introducing z and remains so; it need not be duplicated beneath the constraints.

{\small Evidence: model-analysis.tex, joint continuous problem.}

\subsection*{R16: Continuous measure primal}

The main gives the measure-relaxation definition and points to the displayed primal D9, its incidence/\allowbreak moment constraints, Dirac embedding and integrated weak-duality argument.

{\small Evidence: model-analysis.tex, full-proofs.tex D9.}

\subsection*{R17: Formal decision language}

IDEAL-MMSE-PAIR is defined before Theorem 1 with the explicit alphabet-index table, unit powers/\allowbreak noise, threshold and malformed-input rejection. Abstract/\allowbreak introduction summarize hardness without an earlier undefined NP-completeness declaration.

{\small Evidence: model-analysis.tex, IDEAL-MMSE-PAIR.}

\subsection*{R18: Bounded-degree 3DM citation}

The existing primary source is appropriate: Dyer--Frieze (1985), Sec.2, statement(2.1), printed p.141 gives the exact perfect planar variant with every element in two or three triples. This locator and restriction are now explicit; no unverified replacement citation is inserted.

{\small Evidence: model-analysis.tex Theorem1, full-proofs.tex D1, discovery/\allowbreak astra-theory.md.}

\subsection*{R19: Polynomial NP verification}

The schedule certificate, constant local rational-factor size, O(K)-bit product and polynomial comparison with 2\textasciicircum{}K are stated. No NP-membership claim for unrestricted continuous thresholds is added.

{\small Evidence: model-analysis.tex, Theorem1 proof sketch.}

\subsection*{R20: Alphabet and separation}

The existing constructive alphabet and positive constant gaps were already in S1/\allowbreak S11. The finite alphabet is now specified deterministically, the input-size bound is explicit and main pointers identify D1--D3. A long vector enumeration is unnecessary for this fixed algorithmically specified alphabet.

{\small Evidence: technical-supplement.tex A1--A3, full-proofs.tex D1--D3.}

\subsection*{R21: Forest representation reference}

The proposed rejection of AIM2008 was not adopted: its Theorem 3.16, p.1639 directly gives the tree-complement antecedent. The main now states that locator and distinguishes the additional rational-unit finite-family construction.

{\small Evidence: model-analysis.tex, after Theorem1, technical-supplement.tex.}

\subsection*{R22: Hessian and sharp boundary}

Full Hessian and determinant formulas were already available as D4--D5 and all coefficients in S2. Precise pointers and boundary details were strengthened; new symbolic/\allowbreak finite-point tests below, at and above the threshold pass.

{\small Evidence: full-proofs.tex D4--D5, technical-supplement.tex S2, checks/\allowbreak theory-regression.json.}

\subsection*{R23: Negative definiteness}

The statement now gives the explicit finite-point matrix inequality over the nonnegative quadrant.

{\small Evidence: model-analysis.tex Theorem2.}

\subsection*{R24: Pricing polynomial definitions}

The proof sketch now points directly to the explicit B,C and Pi formulas in D6--D8.

{\small Evidence: model-analysis.tex Theorem3, full-proofs.tex D6--D8.}

\subsection*{R25: Power-budget collision}

The elimination polynomial is renamed Pi throughout main and supplementary proof; P remains the power budget.

{\small Evidence: model-analysis.tex, full-proofs.tex.}

\subsection*{R26: Coordinates and clipping}

Coordinates j in {x,y}, their interval/\allowbreak price symbols and clip are explicitly defined.

{\small Evidence: model-analysis.tex Theorem3.}

\subsection*{R27: Validated lower value}

L is defined as a rigorously validated feasible lower enclosure before its use in the bound statement.

{\small Evidence: model-analysis.tex Proposition1.}

\subsection*{R28: Fixed-pair gain notation}

The example uses gain symbols for the fixed users, avoiding an unexplained indexing convention.

{\small Evidence: model-analysis.tex Proposition1 discussion.}

\subsection*{R29: Both convex terms}

g\_\allowbreak x and g\_\allowbreak y are explicitly defined and their convexity is attributed to those two terms.

{\small Evidence: model-analysis.tex affine majorants.}

\subsection*{R30: Best-bound termination}

The suggested unqualified fairness lemma is replaced by a direct finite-epsilon proof with evaluated feasible lower corners, finite assignment roots and longest-edge subdivision. Experimental finite caps and fixed precision remain outside that unrestricted convergence statement.

{\small Evidence: model-analysis.tex Proposition2, full-proofs.tex final convergence subsection.}

\subsection*{R31: Retained upper bound}

U is defined as the minimum retained validated global upper bound, including fallback and available complete bounds.

{\small Evidence: model-analysis.tex bound/\allowbreak checker definitions.}

\subsection*{R32: Checker reproducibility}

S12 now supplies canonical serialization, exact stored-value reconstruction, the independent 100-bit/\allowbreak 40-term log enclosure M5, rejection/\allowbreak deadline behavior and the actual certificates replay command. The finite-candidate oracle remains a separate algorithm.

{\small Evidence: checked-procedures.tex M5, revision/\allowbreak CERTIFICATE\_\allowbreak API\_\allowbreak AND\_\allowbreak TRUST.md.}

\subsection*{R33: DeepMIMO identity and axes}

Version 0.2.9, X--Y--Z axis order, full configs, installed-source and input hashes are archived. Historical VCS commit and semantic scenario tags were not recorded; exact hashes replace invented labels, with the limitation explicit.

{\small Evidence: revision-details.tex S15, revision/\allowbreak PROVENANCE.json, GENERATOR\_\allowbreak CONFIGS.json.}

\subsection*{R34: Normalization and bandwidth}

The actual no\_\allowbreak LPF code uses sqrt(path\_\allowbreak power/\allowbreak NFFT); sqrt512 removes that path-amplitude factor, not an inferred IFFT convention. The 32 stored tones occupy 3.125 MHz in aggregate, while evaluated 4/\allowbreak 8-tone cases occupy 0.390625/\allowbreak 0.78125 MHz.

{\small Evidence: revision-details.tex R1, evaluation.tex.}

\subsection*{R35: Case provenance}

A 32-case manifest records case IDs, seeds, generator row ranges, selected pool offsets, FFT indices, case/\allowbreak channel/\allowbreak pool hashes. Offsets are not mislabeled as physical row IDs.

{\small Evidence: revision/\allowbreak CASE\_\allowbreak MANIFEST.json, GENERATOR\_\allowbreak CONFIGS.json.}

\subsection*{R36: Matching--WMMSE details}

Already available equations are now numbered M1--M4 and supplemented by Algorithm S1, exact seeds, all starts, assignment scores, bisection/\allowbreak repair, stopping and exchange limits. The retained equation-to-code map identifies implementation functions.

{\small Evidence: checked-procedures.tex, revision/\allowbreak MATCHING\_\allowbreak EQUATION\_\allowbreak CODE\_\allowbreak MAP.md.}

\subsection*{R37: Allowance comparison}

The comparator caption explicitly distinguishes 18 s at K=8 from 28 s at K=16 and states equality only within each size.

{\small Evidence: evaluation.tex comparator table.}

\subsection*{R38: Historical environment}

The recorded CPU, 20/\allowbreak 28 core counts, RAM, OS build, Python/\allowbreak libraries, distribution hashes and single-thread policy are archived. Historical BLAS build and an independently recorded HiGHS version remain explicitly unknown. No current installation is misrepresented as historical evidence.

{\small Evidence: revision-details.tex, revision/\allowbreak PROVENANCE.json.}

\subsection*{R39: Statistical citations}

Primary Clopper--Pearson and Efron references and the official NIST general Bonferroni construction are cited at the method introductions.

{\small Evidence: evaluation.tex, references.bib, revision/\allowbreak REFERENCE\_\allowbreak ADDITIONS.json.}

\subsection*{R40: Estimand and intervals}

The main names the equal-scene discordance estimand; S15 gives R2--R3, endpoint arithmetic, missing-case rules and exact paired-bootstrap seed/\allowbreak draw order. The original inference is unchanged.

{\small Evidence: evaluation.tex, revision-details.tex R2--R3.}

\subsection*{R41: Execution counts}

The paper explicitly decomposes 128 executions into 96 main, 24 scale and 8 analytic reliability runs, separate from eight original mechanism cases.

{\small Evidence: evaluation-results.tex, revision-details.tex.}

\subsection*{R42: Receipt versus check}

Receipt publication and mathematical reconstruction are distinguished; 417 is a cache-hit count after per-receipt integrity/\allowbreak timing verification, not an unexplained total of checks.

{\small Evidence: evaluation-results.tex, checked-procedures.tex, revision-details.tex.}

\subsection*{R43: Jain attribution}

The original 1984 DEC TR-301 source is cited; its 1998 archival posting is distinguished from the report date.

{\small Evidence: evaluation-results.tex, references.bib Jain1984.}

\subsection*{R44: Minimum-user rate evidence}

The main wording is narrowed to discussed comparisons, and a supplementary table reports mean minimum-user SE and Jain index for both scenarios and all four main arms.

{\small Evidence: fairness-summary.tex, revision-details.tex.}

\subsection*{R45: Relative-regret meaning}

The main explicitly derives (V\_\allowbreak C-L)/\allowbreak V\_\allowbreak C <= (U-L)/\allowbreak L for L>0, explaining that the plotted percentage is a conservative bound rather than measured loss.

{\small Evidence: evaluation-results.tex, sixteen-user results.}

\subsection*{R46: CSI/\allowbreak GPU evidence}

These adverse results were already documented in S5/\allowbreak S6 and remain in the delivered supplement. The discussion preserves them and does not present new acceleration or learned-CSI claims.

{\small Evidence: scope-evidence.tex, supplement S5--S6.}

\subsection*{R47: AI disclosure}

The acknowledgment identifies Codex/\allowbreak GPT-6/\allowbreak Astra, the specific main/\allowbreak supplement sections and types of assistance. It does not invent independent human verification or final coauthor approval.

{\small Evidence: acknowledgments.tex, revision/\allowbreak REFERENCE\_\allowbreak ADDITIONS.json.}

\subsection*{R48: Artifact identity and public release}

PDP2-CR030-R1 has complete scientific-file hashes plus archive/\allowbreak delivery manifests. This supplies a fixed local artifact identity. Public visibility/\allowbreak release verification remains an explicitly separate blocker; no DOI, upload or public-access claim is fabricated.

{\small Evidence: acknowledgments.tex, revision/\allowbreak REPRODUCTION\_\allowbreak MANIFEST.json, DELIVERY\_\allowbreak MANIFEST.json.}

\subsection*{R49: Physical model example}

The AP labels now say evaluated example/\allowbreak M=8; the main retains general receive dimension and identifies the drawing as illustrative geometry.

{\small Evidence: images/\allowbreak system-pairing.pdf/\allowbreak .svg, model-analysis.tex.}

\subsection*{R50: Diagram typography}

Both physical diagrams were regenerated as PDF/\allowbreak SVG with a consistent font and 8-point labels in the native 3.25-inch figure, visually checked at article column width. Long labels were split to avoid overlap.

{\small Evidence: draw\_\allowbreak system.py, images/\allowbreak system-mmse.pdf/\allowbreak .svg.}

\subsection*{R51: Status figure and names}

Figure 3 shows the observed nonzero category with full method names. Its caption explicitly records zero target successes/\allowbreak errors/\allowbreak missing outcomes; no observations are dropped.

{\small Evidence: plot\_\allowbreak figures.py, images/\allowbreak fig01-certification.pdf/\allowbreak .svg.}

\subsection*{R52: Target-line annotation}

The gap CDF directly labels the 0.01-nat target; marker and line styles are redundant for grayscale.

{\small Evidence: images/\allowbreak fig02-gap-ecdf.pdf/\allowbreak .svg.}

\subsection*{R53: Zero rate differences}

Hollow triangles drawn over the zero reference line expose Full--Grid equality without changing vertical values.

{\small Evidence: images/\allowbreak fig03-rate-differences.pdf/\allowbreak .svg.}

\subsection*{R54: Runtime names and definition}

Figure 6 uses full method names and parent-process wall-clock wording; the caption states startup/\allowbreak search/\allowbreak check/\allowbreak persistence accounting.

{\small Evidence: images/\allowbreak fig04-charged-duration.pdf/\allowbreak .svg, evaluation-results.tex.}

\subsection*{R55: Floor censoring}

Open downward triangles mark sub-floor enclosures; horizontal offsets distinguish coincident cases, with no fabricated zero or altered gap value.

{\small Evidence: images/\allowbreak fig05-curvature-mechanism.pdf/\allowbreak .svg.}

\subsection*{R56: Fairness zoom}

The O1 panel magnifies 0.98--1; all observations are asserted within view. The complete supplementary fairness/\allowbreak minimum-rate table preserves the numeric context.

{\small Evidence: images/\allowbreak fig06-fairness.pdf/\allowbreak .svg, fairness-summary.tex.}

\subsection*{R57: Equal scale rates}

Open/\allowbreak filled markers and explicitly disclosed small horizontal offsets show Grid/\allowbreak Grid+pricing equality without changing achieved rates.

{\small Evidence: images/\allowbreak fig07-scale-rates.pdf/\allowbreak .svg.}

\subsection*{R58: Scale bar values}

Bars carry one-decimal values and distinct hatches; captions explain relative-regret bounds and complete pricing coverage.

{\small Evidence: images/\allowbreak fig08-scale-certificates.pdf/\allowbreak .svg.}

\subsection*{Additional checklist and presentation conventions}

The compact notation table is in S15, while essential symbols remain defined at first use in the main. Table I is titled ``Comparison with the closest prior work,'' with scope limitations in a note. The reference-algorithm summary is retained in Section V; exact Matching--WMMSE pseudocode is Algorithm S1. Main performance plots use Grid, Curvature, Full, Matching--WMMSE and Grid+pricing consistently. Mechanism controls retain the distinct Reference arm because it is a different experiment, not a synonym for Grid. The historical four-AP supplement is explicitly separate and preserves historical method meanings.

The ten current main figures are supplied as vector PDF and SVG, with consistent typography and grayscale-redundant markers, lines or hatches. Horizontal display offsets never change measured rates/\allowbreak gaps. SE summaries use three/\allowbreak four decimals, gap means three decimals and reported confidence limits four decimals; the 0.01-nat SE conversions are rounded to 0.00361/\allowbreak 0.00180 bit/\allowbreak s/\allowbreak Hz. Case/\allowbreak config/\allowbreak serialization constants and hashes retain exact digits where reproducibility requires them.

The current IEEE/\allowbreak TCOM author instructions were checked from the official policy page. The applicable initial limit remains 13 pages, two columns, 10-point body and one-inch margins; no 16-page journal-revision allowance is claimed. AI assistance is disclosed by system, section and role. The new acknowledgment deliberately avoids the review's suggested claim that the human authors independently verified every result, because that confirmation has not been supplied.

Existing bibliography entries used in the supplement are retained even when absent from the main. The two author-paper and three closest-external-paper exports remain separate and all five are cited substantively. The corresponding-author identity and cover-letter factual scope are preserved; no submission, message, public upload or payment occurred.

\end{document}
